\section{Further research questions}
\label{dhr:sec:research}

\subsection{Higher resonances: several divergent moment entries}
The first resonance $b-a-c=2r-2$ is special because only one particle
can escape at logarithmic order while the lower ensemble retains a
finite first moment. At $b-a-c=2r-2-M$ with $M\geq1$, a larger
block of moments is divergent and the Gamma product can have higher
pole order. Determine a canonical largest-index or multi-index
subtraction, including all overlap terms, and compare its hard-cutoff
polynomial in $\log N$ with its spectral Laurent coefficients.
A successful extension must prove normal convergence after the full
subtraction, not merely expand the Gamma product.

A precise initial target is the next resonance for $r=2$, where one
can explicitly compare the two divergent moments in the cutoff
Hankel determinant with the two relevant gamma poles. This should
clarify which derivatives of residue amplitudes enter beyond the
first-order rule \eqref{dhr:eq:spectral-vs-cutoff}.

\subsection{Repeated rational poles and the rank of harmonic data}
The simple-pole hypothesis in Theorem~\ref{dhr:thm:rank-one} is
essential to the proof. With repeated poles, moments contain
$H_{N+1}^{(2)},H_{N+1}^{(3)},\ldots$ as well as $H_{N+1}$.
Describe the matrix subspace spanned by these harmonic coefficients.
Can its rank and support give a sharp \emph{formal} upper bound on
the polylogarithmic depth of the cutoff generator?

One concrete problem is to classify all rational weights for which
this harmonic matrix pencil has rank one despite repeated poles.
Such a classification could extend the current depth-two theorem
without relying on generic determinant-degree estimates.

\subsection{Closed all-rank coefficients and discrete orthogonal polynomials}
The determinant formulas for $P_r(N)$ and $Q_r(N)$ are exact at every
rank, but expression growth is substantial. Find product formulas,
short recurrences, or a discrete-orthogonal-polynomial interpretation
for the rational coefficients in the canonical case. The factors of
$P_r$ displayed at ranks two through four suggest more structure
than the determinant proof alone explains.

A useful target is a recurrence in $r$ with rational coefficients in
$N$ that generates both $P_r$ and $Q_r$ and preserves the initial
vanishing $S_{r,N}=0$ for $N<r-1$. A companion recurrence for
$I_{r,q}$ could explain systematic weight drops in
\eqref{dhr:eq:I31}--\eqref{dhr:eq:I41}.

\subsection{Mixed spectral and shape derivatives}
Theorem~\ref{dhr:thm:resonance} gives joint holomorphy in $a,c,u$.
One can therefore differentiate in $a,c$ as well as $u$, provided
one keeps track of the moving relation $b_0=a+c+2r-2$.
Develop explicit harmonic/polygamma formulas for these tangential
shape derivatives and compare them with derivatives taken at fixed
$b$. The distinction changes the counterterm and is likely important
for integration with the existing mixed-Dougall calculus.

A targeted problem is to derive a compact three-parameter Bell table
whose rows express mixed generalized harmonic sums in terms of
$\gamma_k(\alpha)$ and polygamma values at $a+r-1,c+r-1$.
The analytic permission to differentiate is proved here; the most
useful arithmetic basis for the resulting rows is not selected.

\subsection{Rational grids beyond a common translate}
The root filters in \eqref{dhr:eq:single-filter} and
\eqref{dhr:eq:double-filter} handle one common residue class.
For rational poles from several classes modulo a common denominator,
the moment matrix can involve several independent shifted harmonic
families. Determine whether its harmonic rank, rather than the number
of summation variables, controls colored depth.

The first test case should use two residue classes with separate
harmonic coefficients and an explicitly computed matrix pencil.
Any resulting endpoint reduction must state its cyclotomic level and
its precise relation system; ordinary-zeta closure should not be
assumed from the common-integer-grid case.

\subsection{Changing the Vandermonde exponent}
The square $\Delta^2$ is what permits Cauchy--Binet and a moment
determinant. For $\Delta^{2\kappa}$ with integer $\kappa>1$, one
can still expand finitely in monomials, but the single-entry argument
no longer applies automatically. Determine the first resonant decay,
its subtraction amplitude, and a useful depth bound. A comparison
with known discrete Selberg-type sums should precede any novelty
claim in that direction.

\subsection{Functional certificates and arithmetic conjectures}
The current identities live in a tightly specified depth-two family.
It is reasonable to ask whether translated or colored variants supply
new functional relations useful for the manuscript's arithmetic
conjectures. There is no result here showing that they settle $S_6$
or $S_8$. A productive target would be an exact functional certificate
whose evaluation gives a previously missing relation, not another
high-precision residual or an unsupported claim of period dependence.

\subsection{Formalization and literature reconciliation}
The rational determinant lemma, finite Cauchy--Binet identity, residue
sum criterion, and terminating transform are natural first targets
for formalization. The analytic normal-convergence argument should
be formalized separately, with its compact-uniform first-moment
estimate and Cauchy derivative transfer made explicit.

In parallel, a systematic literature audit should identify which
forms of the Gauss--Vandermonde product and height-one generating
functions already occur in discrete determinant and Euler-sum
literature, then isolate the exact scope of the resonant cutoff and
rank-one harmonic statements. The present report supplies proofs and
reproducible identities; it does not turn an incomplete priority
search into a novelty theorem.
